% ============================================================================
% Model 4: 薛定谔化方法（面向 PDE 数值计算研究生的 beamer 课件）
% 编译方式: xelatex (两遍)
%   基于完善后的笔记 Notes/Model4_Schrodingerization 与暑期学校 day18/19 课件
%   与 Model 1/2/3 同一风格；本模块完成"从量子算法到 PDE 求解"的最后一块拼图
% ============================================================================
\documentclass[aspectratio=169,9pt]{ctexbeamer}
\usetheme{default}
\usefonttheme{professionalfonts}
\usepackage{amsmath,amssymb,mathtools,bm,braket}
\usepackage{booktabs,array,multirow}
\usepackage{tikz}
\usetikzlibrary{arrows.meta,positioning,calc,decorations.pathreplacing,shapes.geometric,fit,backgrounds}
\usepackage{quantikz}
\usepackage[most]{tcolorbox}
\usepackage{physics}
\usepackage{hyperref}

% --- 配色：与 Model 1/2/3 一致 ---
\definecolor{navy}{RGB}{0,0,145}
\definecolor{deepred}{RGB}{150,0,0}
\definecolor{softblue}{RGB}{238,242,255}
\definecolor{softred}{RGB}{255,238,238}
\definecolor{softgreen}{RGB}{238,250,238}
\definecolor{softyellow}{RGB}{255,248,225}
\definecolor{darkgreen}{RGB}{30,105,50}
\definecolor{graytext}{RGB}{70,70,70}

\setbeamercolor{frametitle}{fg=white,bg=navy}
\setbeamerfont{frametitle}{series=\bfseries,size=\large}
\setbeamercolor{title}{fg=black,bg=white}
\setbeamercolor{structure}{fg=navy}
\setbeamercolor{block title}{fg=white,bg=navy}
\setbeamercolor{block body}{fg=black,bg=softblue}
\setbeamercolor{alerted text}{fg=deepred}
\setbeamertemplate{navigation symbols}{}
\setbeamertemplate{itemize item}{\color{navy}$\blacktriangleright$}
\setbeamertemplate{itemize subitem}{\color{navy}$\triangleright$}
\setbeamertemplate{enumerate item}{\color{navy}\insertenumlabel.}
\setbeamertemplate{frametitle}{%
  \nointerlineskip
  \begin{beamercolorbox}[wd=\paperwidth,ht=2.9ex,dp=0.9ex,leftskip=1.0em]{frametitle}
    \insertframetitle
  \end{beamercolorbox}}
\setbeamertemplate{footline}{%
  \hfill\scriptsize\color{graytext}\insertframenumber{} / \inserttotalframenumber\hspace{1.5em}\vspace{0.7em}}

\setbeamersize{text margin left=0.9em,text margin right=0.9em}
\setbeamertemplate{itemize/enumerate body begin}{\small}
\setbeamertemplate{itemize/enumerate subbody begin}{\footnotesize}
\setlength{\abovedisplayskip}{4pt}
\setlength{\belowdisplayskip}{4pt}
\setlength{\abovedisplayshortskip}{3pt}
\setlength{\belowdisplayshortskip}{3pt}

\newtcolorbox{keybox}[1]{colback=softblue,colframe=navy,title={#1},fonttitle=\bfseries,boxrule=0.5pt,arc=1mm,left=1mm,right=1mm,top=0.6mm,bottom=0.6mm,boxsep=0.5mm}
\newtcolorbox{warnbox}[1]{colback=softred,colframe=deepred,title={#1},fonttitle=\bfseries,boxrule=0.5pt,arc=1mm,left=1mm,right=1mm,top=0.6mm,bottom=0.6mm,boxsep=0.5mm}
\newtcolorbox{exbox}[1]{colback=softyellow,colframe=darkgreen,title={#1},fonttitle=\bfseries,boxrule=0.5pt,arc=1mm,left=1mm,right=1mm,top=0.6mm,bottom=0.6mm,boxsep=0.5mm}

\newtheorem*{Def}{定义}
\newtheorem*{Thm}{定理}
\newtheorem*{Lem}{引理}
\newtheorem*{Cor}{推论}
\newtheorem*{Prop}{命题}

\newcommand{\C}{\mathbb{C}}
\newcommand{\R}{\mathbb{R}}
\newcommand{\N}{\mathbb{N}}
\newcommand{\eps}{\varepsilon}
\newcommand{\poly}{\operatorname{poly}}
\newcommand{\ii}{\mathrm{i}}
\newcommand{\e}{\mathrm{e}}
\newcommand{\smallcite}[1]{\par\vspace{1.2mm}{\scriptsize\color{graytext}#1}}

\title{Schrödingerisation\\[2mm] \large 薛定谔化方法：从经典 PDE 到酉量子模拟}
\author{基于 Model 4 笔记与暑期学校 day18/19 课件整理}
\date{\today}

\begin{document}

% ============================================================================
% 封面与课程定位
% ============================================================================
\begin{frame}[plain]
\vspace{0.9cm}
{\LARGE\bfseries Schrödingerisation}\par\vspace{3mm}
{\large 薛定谔化方法：从经典 PDE 到酉量子模拟}\par\vspace{6mm}
\begin{center}
\begin{tikzpicture}[node distance=4mm, every node/.style={font=\small,align=center}]
\node[draw,rounded corners,fill=white,minimum width=2.0cm,minimum height=0.8cm] (pde) {线性 PDE\\ $u_t=Lu$};
\node[draw,rounded corners,fill=softblue,minimum width=2.2cm,minimum height=0.8cm,right=of pde] (schr) {薛定谔化\\ warped phase};
\node[draw,rounded corners,fill=softgreen,minimum width=2.0cm,minimum height=0.8cm,right=of schr] (disc) {离散化\\ $x,\xi$ 网格};
\node[draw,rounded corners,fill=softyellow,minimum width=2.2cm,minimum height=0.8cm,right=of disc] (be) {Block-encoding\\ $H_{\mathrm{eff}}$};
\node[draw,rounded corners,fill=softred,minimum width=2.2cm,minimum height=0.8cm,right=of be] (hs) {酉演化\\ $e^{-iH_{\mathrm{eff}}T}$};
\node[draw,rounded corners,fill=white,minimum width=2.0cm,minimum height=0.8cm,right=of hs] (out) {测量读出\\ $u(x,T)$};
\draw[-Latex,thick] (pde)--(schr); \draw[-Latex,thick] (schr)--(disc);
\draw[-Latex,thick] (disc)--(be); \draw[-Latex,thick] (be)--(hs); \draw[-Latex,thick] (hs)--(out);
\end{tikzpicture}
\end{center}
\vfill
{\small 把一般线性非酉 PDE/ODE \textbf{整体}转化为薛定谔方程，只做\textbf{一次}哈密顿量模拟。
本模块是系列课程的收官：Model 1 语言 $\to$ Model 2 构件 $\to$ Model 3 统一框架 $\to$ Model 4 科学计算。}
\smallcite{素材：Model 4 笔记 \emph{Schrodingerization}；湘潭大学暑期学校 day18/19 课件。}
\end{frame}

\begin{frame}{从线性代数到微分方程：两条主线}
前三个模块建立了"语言 $\to$ 构件 $\to$ 统一框架"的链条；本模块把构件用于\textbf{科学计算}。
\begin{columns}[T,totalwidth=\textwidth]
\column{0.5\textwidth}
\begin{keybox}{主线 1：HHL 线性方程组}
\[
A x = b \;\longrightarrow\; \ket x = \frac{A^{-1}\ket b}{\norm{A^{-1}\ket b}}
\]
\begin{itemize}
\item 路线：微分方程 $\to$ 时间离散 $\to$ 线性方程组 $\to$ 量子线路；
\item 构件：QPE + 受控旋转（或 QSVT 直接求逆）。
\end{itemize}
\end{keybox}
\column{0.5\textwidth}
\begin{keybox}{主线 2：薛定谔化}
\[
u_t = Lu \;\longrightarrow\; i\partial_t\hat v = (\xi H_1 + H_2)\hat v
\]
\begin{itemize}
\item 路线：PDE $\to$ warped phase transformation $\to$ 一次哈密顿量模拟；
\item 构件：block-encoding + Hamiltonian 模拟 + 振幅放大。
\end{itemize}
\end{keybox}
\end{columns}
\begin{warnbox}{共享构件}
两条主线共享同一套底层构件（block-encoding、哈密顿量模拟、振幅放大）——
所以 §3 先对哈密顿量模拟做统一综述，作为两者的共同基础。
\end{warnbox}
\smallcite{Model 4 笔记第 1 节。}
\end{frame}

\begin{frame}{本模块知识地图}
\begin{columns}[T,totalwidth=\textwidth]
\column{0.5\textwidth}
\begin{keybox}{§2 HHL 算法}
\begin{itemize}
\item QLSP 问题、厄米膨胀与归一化；
\item QPE + 受控旋转 + 逆 QPE；
\item 复杂度：$O(\kappa^2/\eps)$，含 AA；
\item QSVT-HHL 与经典方法对比。
\end{itemize}
\end{keybox}
\begin{keybox}{§3 哈密顿量模拟综述}
\begin{itemize}
\item 任务模型、三种输入模型、无快速前推下界；
\item Trotter/Suzuki（交换子界、高阶、随机化）；
\item Taylor/LCU（BCCKS）；Qubitization/QSVT（最优）；
\item 含时：相互作用绘景。
\end{itemize}
\end{keybox}
\column{0.5\textwidth}
\begin{keybox}{§4 薛定谔化（核心）}
\begin{itemize}
\item Warped phase transformation（热方程型、一般型）；
\item 恢复定理（有限传播速度论证）；
\item Smooth initialization（$\mu_{\max}=O(\log 1/\eps)$）；
\item 源项、时间依赖、复杂度与量子优势、病态 PDE。
\end{itemize}
\end{keybox}
\begin{keybox}{§5 酉化方法族}
薛定谔化 / LCHS / Transmutation / Moment-matching dilation 的统一视角与选择。
\end{keybox}
\end{columns}
\smallcite{Model 4 笔记各节。}
\end{frame}

% ============================================================================
\section{HHL 算法：求解线性方程组}
\begin{frame}[plain]
\centering\vspace{2.0cm}
{\Huge\bfseries\color{navy} HHL 算法：求解线性方程组}\par\vspace{4mm}
{\large HHL: Quantum Linear System Solver}
\end{frame}

\begin{frame}{QLSP 问题设定与厄米假设}
\begin{Def}[\textbf{QLSP}]
给定 $n$ 比特矩阵 $A\in\C^{N\times N}$（$N=2^n$）与归一化右端 $\ket b=U_b\ket{0^n}$，求解 $Ax=b$。
目标是制备归一化解态
\[
\ket x = \frac{A^{-1}\ket b}{\norm{A^{-1}\ket b}},
\]
并使其与真实解满足 $\norm{\ket x - \ket{\bar x}}\le\eps$。
\end{Def}
\begin{keybox}{厄米假设与膨胀（不失一般性）}
非厄米 $A$ 用\textbf{厄米膨胀}
$\bar A=\begin{pmatrix}0&A\\A^\dagger&0\end{pmatrix}
=\ket1\bra0\otimes A+\ket0\bra1\otimes A^\dagger$（$\bar A\ket0\ket x=\ket1\ket b$）；
再归一化 $A\to A/\norm{A}$ 使 $\norm{A}=1$，于是条件数
\[
\kappa := \norm{A}\,\norm{A^{-1}} = \norm{A^{-1}} = 1/\lambda_0
\qquad (0<\lambda_0\le\dots\le\lambda_{N-1}=1).
\]
归一化只改变解的整体尺度，归一化解态不变。
\end{keybox}
\smallcite{Model 4 笔记 §2.1。}
\end{frame}

\begin{frame}{HHL 三步骤}
设 $A=\sum_j\lambda_j\ket{v_j}\bra{v_j}$，$\ket b=\sum_j\beta_j\ket{v_j}$，则
$A^{-1}\ket b=\sum_j\beta_j\lambda_j^{-1}\ket{v_j}$ —— 只需"读出"本征值并乘 $\lambda_j^{-1}$（QPE 的强项）。
\begin{Def}[\textbf{HHL}（$U=e^{i2\pi A}$ 由哈密顿量模拟实现）]
\begin{enumerate}
\item \textbf{QPE}：$U_{\mathrm{QPE}}\ket{0^d}\ket b=\sum_j\beta_j\ket{\lambda_j}\ket{v_j}$（本征值进相位寄存器）；
\item \textbf{受控旋转}：$U_{\mathrm{CR}}\ket0\ket{\lambda_j}=\big(\sqrt{1-C^2/\lambda_j^2}\ket0+\tfrac{C}{\lambda_j}\ket1\big)\ket{\lambda_j}$（$0<C\le\lambda_0$）；
\item \textbf{逆 QPE（uncomputation）}：$U_{\mathrm{QPE}}^\dagger$ 复位相位寄存器为可复用工作寄存器。
\end{enumerate}
\end{Def}
\begin{center}
\begin{quantikz}[row sep=0.2cm,column sep=0.26cm]
\lstick{$\ket0$} & \qw & \qw & \qw & \qw & \qw & \meter{} & \rstick{0/1} \\
\lstick{$\ket{0^d}$} & \gate{U_{\mathrm{QPE}}} & \gate{U_{\mathrm{CR}}} & \gate{U_{\mathrm{QPE}}^\dagger} & \qw & \qw & \qw & \qw \\
\lstick{$\ket b$} & \qw & \qw & \qw & \qw & \qw & \qw & \rstick{$\ket x$}
\end{quantikz}
\end{center}
\begin{warnbox}{电路约定}
$U_{\mathrm{QPE}}$ 同时（受控地）作用于系统寄存器（系统线"穿过"方框）；
$U_{\mathrm{CR}}$ 以相位寄存器（内容 $\lambda_j$）为控制。
\end{warnbox}
\smallcite{Model 4 笔记 §2.2。}
\end{frame}

\begin{frame}{HHL 的输出与范数恢复}
\begin{Thm}[\textbf{HHL 的输出}]
上述电路实现
\[
U_{\mathrm{HHL}}\ket0\ket b=\sum_j\beta_j\Big(\sqrt{1-C^2/\lambda_j^2}\ket0+\frac{C}{\lambda_j}\ket1\Big)\ket{v_j}.
\]
测量信号比特得到 $1$ 时，系统寄存器处于（归一化的）
$\ket x\propto\sum_j\beta_j\frac{C}{\lambda_j}\ket{v_j}=C A^{-1}\ket b$
—— 即所求近似解（常数 $C$ 在归一化后消失）。
\end{Thm}
\begin{keybox}{恢复解的范数}
成功概率 $p(1)=\norm{CA^{-1}\ket b}^2$：通过估计 $p(1)$（重复运行电路）可以恢复
$\norm{A^{-1}\ket b}$ —— 这对需要解向量范数的下游任务（如 PDE 后处理）很重要。
\end{keybox}
\begin{exbox}{受控旋转的实现}
取 $\theta_j=\arcsin(C/\lambda_j)$，用经典算术电路实现 $U_{\mathrm{angle}}:\ket{\lambda_j}\mapsto\ket{\theta_j}$
（$O(\mathrm{poly}(d'))$ 门），施加 $R_y(2\theta_j)$ 后再 uncomputation；
相位寄存器位数 $d'=O(\log(\kappa/\eps))$ 即可保证精度（误差线性传播，Model 1）。
\end{exbox}
\smallcite{Model 4 笔记 §2.2--2.3。}
\end{frame}

\begin{frame}{HHL 的复杂度}
\begin{Thm}[\textbf{查询复杂度}]
设 $\norm{A}=1$，$\lambda_0=1/\kappa$，取 $C=\lambda_0=1/\kappa$：
\begin{enumerate}
\item \textbf{成功概率}：$p(1)=\norm{CA^{-1}\ket b}^2\ge C^2/\norm{A}^2=C^2=\Omega(\kappa^{-2})$
（由 $\norm{b}=\norm{AA^{-1}b}\le\norm{A}\norm{A^{-1}b}$，$\norm{A}=1$ 得 $\norm{A^{-1}b}\ge1$）；
\item \textbf{精度}：QPE 需对本征值做到加性精度 $\eps/\kappa$（$\delta_j$ 使 $\lambda_j^{-1}$ 有相对误差
$O(\delta_j/\lambda_j)$），单次运行对 $U=e^{i2\pi A}$ 查询 $O(\kappa/\eps)$ 次；
\item \textbf{总复杂度}：直接重复 $O(\kappa^2)$ 次得 $O(\kappa^3/\eps)$；振幅放大把重复次数降为
$O(\kappa)$，得
\[
\boxed{\ \text{HHL 查询复杂度} = O(\kappa^2/\eps)\ (\text{含 AA})},\qquad
O(\kappa/\eps)\ (\text{最佳情形 } \norm{A^{-1}\ket b}=\Omega(1)).
\]
\end{enumerate}
\end{Thm}
\begin{keybox}{AA 的适用性}
好态由单个信号比特 $\ket1$ 标志，$R_{\mathrm{good}}=Z\otimes I$；关于初态的反射由
$U_0=U_{\mathrm{HHL}}(I\otimes U_b)$ 构造。$O(\kappa^3/\eps)\to O(\kappa^2/\eps)$ 正是
AA 把 $O(\kappa^2)$ 次重复降为 $O(\kappa)$ 次的结果。
\end{keybox}
\smallcite{Model 4 笔记 §2.4。}
\end{frame}

\begin{frame}{HHL 的现代替代：QSVT-HHL}
\begin{keybox}{QPE 路线 vs QSVT 路线}
\begin{itemize}
\item \textbf{多项式度}：在 $[-1,-1/\kappa]\cup[1/\kappa,1]$ 上逼近 $1/(\kappa x)$ 的奇多项式度数
$O(\kappa\log(\kappa/\eps))$；
\item \textbf{查询复杂度}：$O\big(\frac{\kappa^2}{\zeta}\log(\kappa/\eps)\big)$（$\zeta=\norm{A^{-1}\ket b}\ge1$；
$\ket b$ 与最小奇异方向有 $\Omega(1)$ 重叠时降为 $O(\kappa\log(\kappa/\eps))$）；
\item \textbf{精度依赖}：$\log(1/\eps)$ 代替 $1/\eps$；\textbf{无需厄米/膨胀}（直接奇异值变换）。
\end{itemize}
\end{keybox}
\begin{center}
\footnotesize
\begin{tabular}{ccc}
\toprule
方法 & 查询复杂度（主导项） & 条件\\
\midrule
HHL（QPE） & $\tilde O(\kappa^2/\eps)$ & 需厄米 + 哈密顿量模拟\\
QSVT-HHL & $\tilde O(\kappa^2\log(\kappa/\eps))$ & 只需 block-encoding\\
\bottomrule
\end{tabular}
\end{center}
\begin{warnbox}{与经典解法的对比（诚实边界）}
HHL 对 $N$ 是指数改进（$\poly\log N$ vs $O(N)$--$O(N^3)$），但三个前提：
(1) $e^{i2\pi A}$ 或 $U_A$ 须 $\poly(n)$ 门实现（稀疏保证，稠密不成立）；
(2) 输出是量子态，完整读出需 $\Omega(N)$ 次采样，只在"子程序"或"低维观测量"时有优势；
(3) $\kappa$ 依赖（线性/二次）远差于经典迭代法，病态时可能抵消规模优势。
\end{warnbox}
\smallcite{Model 4 笔记 §2.5--2.6。}
\end{frame}

% ============================================================================
\section{哈密顿量模拟综述}
\begin{frame}[plain]
\centering\vspace{2.0cm}
{\Huge\bfseries\color{navy} 哈密顿量模拟综述}\par\vspace{4mm}
{\large Hamiltonian Simulation: Review for the Common Foundation}
\end{frame}

\begin{frame}{任务模型与输入模型}
\begin{Def}[\textbf{哈密顿量模拟}]
给定 $H=H^\dagger$，构造 $\widetilde U(t)$ 使 $\norm{\widetilde U(t)-e^{-iHt}}\le\eps$
（算子范数意义、与初态无关的态到态精度）。
\end{Def}
\begin{keybox}{任务模型（day17 约定）}
计入：$U_H$（或稀疏 oracle）查询次数、逻辑门数与线路深度、postselection/振幅放大重复；
不计入：初态制备、完整经典读出、物理层纠错。区分"查询复杂度"与"门复杂度"：
前者隐藏 $U_H$ 的实现代价，两者可能相差巨大。
\end{keybox}
\begin{keybox}{三种输入模型}
\begin{itemize}
\item \textbf{稀疏 oracle}：$s$-稀疏 $H$，行/列索引 oracle $O_c$ 与矩阵元 oracle $O_A$；
\item \textbf{block-encoding}：$(\alpha,a,\eps)$-BE $U_H$，$\alpha\ge\norm{H}$；
\item \textbf{Pauli 分解}：$H=\sum_j c_jP_j$（Trotter 常用）。
\end{itemize}
\end{keybox}
\smallcite{Model 4 笔记 §3.1；day17 课件。}
\end{frame}

\begin{frame}{无快速前推定理}
\begin{Thm}[\textbf{no-fast-forwarding}]
对通过黑盒 oracle 访问的\textbf{一般}稀疏哈密顿量 $H$，模拟 $e^{-iHt}$ 至少需要
\[
\Omega(t)
\]
次查询（Berry--Ahokas--Cleve--Sanders 2007）。结合精度下界，哈密顿量模拟查询下界为
\[
\Omega\!\Big(\alpha t+\frac{\log(1/\eps)}{\log\log(1/\eps)}\Big),
\]
其中 $\alpha$ 为 block-encoding 归一化因子（稀疏模型取 $\alpha=\norm{H}$）。
\end{Thm}
\begin{keybox}{直觉与例外}
每个 oracle 查询至多贡献 $O(1)$ 的"演化量"，一般矩阵的谱信息无法被压缩。
该下界不排除\textbf{特殊结构}的快速前推（如 $U=R_z(\varphi)$ 时 $U^j$ 代价与 $j$ 无关），
但这类结构在科学计算的一般矩阵中很少见。
\end{keybox}
\smallcite{Model 4 笔记 §3.1；Berry--Ahokas--Cleve--Sanders, CMP 270 (2007).}
\end{frame}

\begin{frame}{Trotter/Suzuki：一阶交换子型误差界}
对 $H=\sum_{j=1}^m H_j$，一阶 Lie--Trotter $S_1(\tau)=\prod_j e^{-i\tau H_j}$。
\begin{Thm}[\textbf{一阶交换子型误差界}]
对两项 $H=H_1+H_2$，由 Duhamel 原理与 $\norm{[e^{-isH_1},H_2]}\le s\norm{[H_1,H_2]}$：
\[
\norm{e^{-i\tau H}-e^{-i\tau H_1}e^{-i\tau H_2}}\le\frac{\tau^2}{2}\norm{[H_1,H_2]}.
\]
$m$ 项情形全局误差（$L$ 步线性累积）：
\[
\norm{e^{-itH}-S_1(t/L)^L}\le\frac{t^2}{2L}\sum_{j<k}\norm{[H_j,H_k]},
\qquad
L=O\!\Big(\frac{t^2\sum_{j<k}\norm{[H_j,H_k]}}{\eps}\Big)\le O\!\Big(\frac{(t\Lambda)^2}{\eps}\Big),
\]
$\Lambda:=\sum_j\norm{H_j}$。交换子界优于朴素范数界：\textbf{不相交项不贡献误差}。
\end{Thm}
\begin{exbox}{横向场 Ising 模型（交换子界更紧）}
$H=-J\sum Z_iZ_{i+1}-g\sum X_i$：范数界 $L=O(n^2t^2/\eps)$；
$[Z_iZ_{i+1},X_k]\ne0$ 仅当 $k=i,i+1$，故 $\|[H_1,H_2]\|=O(n)$，交换子界 $L=O(nt^2/\eps)$。
\end{exbox}
\smallcite{Model 4 笔记 §3.2；Model 2 笔记 §5。}
\end{frame}

\begin{frame}{Trotter/Suzuki：高阶、实现与随机化}
\begin{Thm}[\textbf{高阶 Suzuki 公式}]
二阶（Strang 分裂）局部误差 $O(\tau^3)$；$p$ 阶公式局部误差 $O(\tau^{p+1})$，达到精度 $\eps$ 需
\[
L=O\!\Big(\frac{(t\Lambda)^{1+1/p}}{\eps^{1/p}}\Big)
\]
步，每步调用 $O(m\cdot5^{\lceil p/2\rceil-1})$ 次局部演化（Suzuki 递推，指数个数随阶数指数增长；
实践中 $p=2,4,6$ 最优）。
\end{Thm}
\begin{keybox}{局部演化的实现}
$e^{-i\tau cP}$（$P$ 为 Pauli 串）：CNOT 链计算奇偶性 $+$ 单个 $R_z(2c\tau)$ $+$ CNOT 链还原
（$2(k-1)$ 个 CNOT 与 1 个 $R_z$，见 Model 2 的 $ZZ$ 电路）；对易项之和可直接快速前推；
PDE 离散矩阵的移位/对角算子指数可由移位电路/对角相位门实现。
\end{keybox}
\begin{keybox}{随机化乘积公式（Childs--Ostrander--Su 2019）}
对 $p$ 阶公式做随机化，\textbf{期望}代价改进为
$\widetilde O\big(t^{1+1/(2p)}\eps^{-1/(2p)}\big)$（指数平方根改进），电路仍是乘积公式，适合 NISQ。
\end{keybox}
\smallcite{Model 4 笔记 §3.2；Childs--Ostrander--Su, Quantum 3, 182 (2019).}
\end{frame}

\begin{frame}{Taylor 级数法与 LCU}
\[
e^{-iHt}\approx\sum_{k=0}^{K}\frac{(-it)^k}{k!}H^k,
\qquad \text{每项 } H^k \text{ 用 } k \text{ 次受控 } U_H.
\]
\begin{warnbox}{直接整段截断的两重困难}
\begin{enumerate}
\item $k\lesssim\tau:=\Lambda t$ 时 $\tau^k/k!$ 仍递增（越过 $k\sim e\tau$ 峰值才开始衰减），
截断阶 $K=O\big(\tau+\frac{\log(1/\eps)}{\log\log(1/\eps)}\big)$；
\item LCU 归一化因子 $\sum_{k\le K}\tau^k/k!\approx e^\tau$ \textbf{指数}大，
后选择成功概率 $\sim e^{-2\tau}$，需指数多次振幅放大，不可行。
\end{enumerate}
\end{warnbox}
\begin{keybox}{LCU 引理（Module3 回顾）}
设 $M=\sum_j w_jV_j$（$w_j\ge0$，$V_j$ 酉），$s=\sum_jw_j$。构造
PREPARE $V\ket{0^a}=\sum_j\sqrt{w_j/s}\ket j$ 与 SELECT
$\mathrm{SEL}=\sum_j\ket j\bra j\otimes V_j$，则
$(V^\dagger\otimes I)\mathrm{SEL}(V\otimes I)\ket{0^a}\ket\psi
=\ket{0^a}\otimes\frac{M\ket\psi}{s}+\ket{0^a}^{\perp}\otimes(\cdots)$，
成功概率 $\norm{M\ket\psi}^2/s^2$。
\end{keybox}
\smallcite{Model 4 笔记 §3.3。}
\end{frame}

\begin{frame}{Taylor 级数法：BCCKS 分段复杂度}
BCCKS 把时间分为 $r=\lceil\tau\rceil$ 段（每段 $\Lambda t/r\le1$），每段截断阶
$K=O\big(\frac{\log(\tau/\eps)}{\log\log(\tau/\eps)}\big)$，使每段误差 $\le\eps/r$、归一化因子 $O(1)$，
再用\textbf{非感知振幅放大}把成功概率提升到常数：
\begin{Thm}[\textbf{截断 Taylor 级数方法的复杂度}]
\[
\boxed{\ O\!\Big(\tau\,\frac{\log(\tau/\eps)}{\log\log(\tau/\eps)}\Big)\ }
\]
对时间 $t$ \textbf{线性}，精度\textbf{指数}依赖（乘性 $\log(1/\eps)$）；
除 $\log\log$ 因子外已达无快速前推下界 $\Omega(t)$。
\end{Thm}
\begin{warnbox}{缺点}
SELECT oracle 需对 $H^k$ 的幂做受控，电路较深；成功概率需振幅放大补偿。
\end{warnbox}
\smallcite{Model 4 笔记 §3.3；Berry et al., PRL 114, 090502 (2015).}
\end{frame}

\begin{frame}{Qubitization 与 QSVT：最优方法}
设 $H$ 有 $(\alpha,m)$-block-encoding $U_H$。定义成功子空间反射
$Z=2\ket{0^m}\bra{0^m}\otimes I-I$ 与 qubitization 迭代算子 $O=U_HZ$：
$(\bra{0^m}\otimes I)O^k(\ket{0^m}\otimes I)=T_k(H/\alpha)$（迭代 $k$ 次给 Chebyshev 多项式，Module3）。
\begin{keybox}{Jacobi--Anger 展开与复杂度}
令 $x=\lambda/\alpha$、$\tau=\alpha t$：$e^{-i\tau x}=J_0(\tau)+2\sum_{k\ge1}(-i)^kJ_k(\tau)T_k(x)$，
故 $e^{-iHt}=J_0(\alpha t)I+2\sum_{k\ge1}(-i)^kJ_k(\alpha t)T_k(H/\alpha)$。
Bessel 系数 $J_k(\tau)\approx(\tau/2)^k/k!$ 在 $k\gtrsim\alpha t$ 后超指数衰减，故
\textbf{深度/查询数} $\boxed{O\big(\alpha t+\frac{\log(1/\eps)}{\log\log(1/\eps)}\big)}$，
匹配下界 $\Omega(\alpha t+\frac{\log(1/\eps)}{\log\log(1/\eps)})$ —— \textbf{最优}（时间与精度加性）。
\end{keybox}
\begin{warnbox}{与 Trotter 的本质区别}
Trotter：时间域\textbf{算子分裂 + 时间步进}（$t^{1+1/p}$ 超线性）；
Qubitization/QSP：\textbf{谱域多项式变换}（对 $t$ 线性、对 $\log(1/\eps)$ 加性）。
\end{warnbox}
\smallcite{Model 4 笔记 §3.4；Low--Chuang, Quantum 3, 163 (2019).}
\end{frame}

\begin{frame}{含时哈密顿量与对比选择}
\begin{keybox}{含时 $H(t)$：相互作用绘景（Low--Wiebe 2018）}
解是时间有序指数 $U(t)=\mathcal T\exp(-i\int_0^t H(s)ds)$。若 $H(t)=A+B(t)$（$A$ 时间无关且易模拟），
在相互作用绘景中代价与 $\int_0^t\norm{B(s)}ds$ 成比例，而非
$t(\norm{A}+\max_s\norm{B(s)})$ —— 小扰动时显著更优（薛定谔化后
$H_{\mathrm{eff}}(\xi)=\xi H_1+H_2$ 中 $\norm{H_2}\ll\norm{H_1}$ 即此类情形）。
\end{keybox}
\begin{center}
\footnotesize
\begin{tabular}{@{}lclcl@{}}
\toprule
方法 & 查询/深度（主导项） & 精度依赖 & 辅助比特 & 适用场景\\
\midrule
Trotter/Suzuki & $O((t\Lambda)^{1+1/p}\eps^{-1/p})$ & 多项式 & $O(1)$ & NISQ, 短时\\
随机化乘积公式 & $\widetilde O(t^{1+1/(2p)}\eps^{-1/(2p)})$ & 期望意义 & $O(1)$ & NISQ, 中等\\
Taylor+LCU & $O(\Lambda t\frac{\log(\Lambda t/\eps)}{\log\log(\Lambda t/\eps)})$ & 指数 & $O(\log K)$ & 容错, 长时\\
Qubitization/QSVT & $O(\alpha t+\frac{\log(1/\eps)}{\log\log(1/\eps)})$ & 指数 & $O(m+1)$ & 容错, 最优\\
\bottomrule
\end{tabular}
\end{center}
\begin{keybox}{薛定谔化管线中的选择}
$H_{\mathrm{eff}}$ 通常是"移位算子 $+$ 对角算子"之和、自然获得 block-encoding，故
QSVT 路线（或结构定制的高阶 Trotter）是首选；NISQ 上 Trotter 仍是务实备选；
若移位与对角部分对易，乘积公式退化为精确快速前推。
\end{keybox}
\smallcite{Model 4 笔记 §3.5--3.6；Low--Wiebe, arXiv:1805.00675.}
\end{frame}

% ============================================================================
\section{薛定谔化方法}
\begin{frame}[plain]
\centering\vspace{2.0cm}
{\Huge\bfseries\color{navy} 薛定谔化方法}\par\vspace{4mm}
{\large Schrödingerisation: The Core}
\end{frame}

\begin{frame}{动机：PDE 的量子模拟}
一般的线性\textbf{发展方程}（热方程、对流扩散、Fokker--Planck、Black--Scholes）空间离散后：
\[
\frac{du}{dt}=Lu,\qquad\text{或更一般地}\quad \frac{du}{dt}=A(t)u+f(t),
\]
其中 $L$ 或 $A$ 一般\textbf{不是厄米矩阵}：
\begin{itemize}
\item 热方程：$L$ 厄米但演化 $e^{-tL}$\textbf{非酉}（耗散）；
\item 对流方程：$A$ 甚至\textbf{反厄米}。
\end{itemize}
\begin{warnbox}{核心困难}
标准哈密顿量模拟只处理 $i\partial_t u=Hu$（$H$ 厄米、演化酉）。
如何把非酉动力学"塞进"酉演化，是 PDE 量子模拟的根本问题。
\end{warnbox}
\begin{keybox}{两条历史路线}
\begin{enumerate}
\item \textbf{线性化 + 线性方程组}：时间离散（隐式欧拉 $(I+\Delta t L)u^{n+1}=u^n$），每步 HHL，共 $M=T/\Delta t$ 步；
\item \textbf{薛定谔化}（Jin--Liu--Yu 2023--2024）：把 $u_t=Lu$ \textbf{整体}化为薛定谔方程，
只做\textbf{一次}哈密顿量模拟。
\end{enumerate}
\end{keybox}
\smallcite{Model 4 笔记 §4.1。}
\end{frame}

\begin{frame}{Warped phase transformation：热方程型}
\begin{Thm}[\textbf{薛定谔化：热方程型}]
设 $u_t=-Lu$，$L=L^\dagger\succeq0$（如 $L=-\Delta$ 加边界条件）。引入辅助变量 $p\ge0$ 与
\textbf{warped phase transformation}
\[
v(t,p)=e^{-p}\,u(t),
\]
则 $\partial_p v=-v$，且 $\partial_t v=-Lv=L\,\partial_p v$。
对 $p$ 做 Fourier 变换（约定 $\hat f(\xi)=\int f(p)e^{-i\xi p}dp$，$\partial_p\leftrightarrow i\xi$）：
\[
i\,\partial_t\hat v(t,\xi)=-\xi L\,\hat v(t,\xi),
\]
即哈密顿量为 $H(\xi)=-\xi L$ 的薛定谔方程 —— 对每个 $\xi\in\R$，$H(\xi)$ 厄米，
演化 $e^{-iH(\xi)t}$ 酉。原解由
\[
u(t)=e^p\,v(t,p)\ (\forall p\ge0),\qquad\text{或}\qquad u(t)=\int_0^\infty v(t,p)\,dp
\]
恢复（$\int_0^\infty e^{-p}dp=1$）。
\end{Thm}
\smallcite{Model 4 笔记定理 4.1。}
\end{frame}

\begin{frame}{为什么叫"warped phase"}
\begin{keybox}{几何直觉}
$v=e^{-p}u$ 在振幅上附加指数衰减权重 $e^{-p}$；$p$ 方向的"相位扭曲"经 Fourier 变换把
\textbf{非酉算子} $L$ 变成\textbf{厄米算子} $-\xi L$，代价是增加一个维度（辅助变量 $\xi$）。
\[
\underbrace{\text{耗散演化 } e^{-tL}}_{\text{非酉}} 
\;\xrightarrow{\text{提升到 }(x,p)\text{ 空间}}\;
\underbrace{\text{保范演化 } e^{-iH(\xi)t}}_{\text{酉}}
\]
直观上：\textbf{耗散被"提升"到高维空间后变为保范演化，通过额外维度记录耗散的历史}。
真正的非酉性没有在线路中出现；线路只做 unitary simulation，非酉效果通过辅助变量 $p$ 的投影恢复。
\end{keybox}
\smallcite{Model 4 笔记 §4.2 注；day18 课件。}
\end{frame}

\begin{frame}{一般情形的薛定谔化}
\begin{Thm}[\textbf{一般线性系统}]
设 $u_t=A(t)u+f(t)$，把 $A$ 分解为厄米部分与反厄米部分
\[
A=H_1-iH_2,\qquad H_1=\frac{A+A^\dagger}{2},\qquad H_2=\frac{i(A-A^\dagger)}{2}
\]
（$H_1,H_2$ 均厄米；$H_1$ 对应增长/耗散，$H_2$ 对应相位旋转；\textbf{不要求} $-H_1\succeq0$）。
做 $v=e^{-p}u$（$\partial_p v=-v$）。关键：把 $H_1$ 项代入 $p$ 导数、$H_2$ 项保留：
\[
\partial_t v=Av=(H_1-iH_2)v=-H_1\partial_p v-iH_2 v.
\]
Fourier 变换（$\partial_p\leftrightarrow i\xi$）：
\[
\partial_t\hat v=-i(\xi H_1+H_2)\hat v
\quad\Longrightarrow\quad
\boxed{\ i\,\partial_t\hat v(t,\xi)=(\xi H_1+H_2)\,\hat v(t,\xi)\ }
\]
即每个频率分量对应一个薛定谔方程，哈密顿量 $H_{\mathrm{eff}}(\xi)=\xi H_1+H_2$ 对每个 $\xi\in\R$ 厄米。
\end{Thm}
\smallcite{Model 4 笔记定理 4.2。}
\end{frame}

\begin{frame}{为什么"只有 $H_1$ 与 $p$ 耦合"是合理的}
\begin{keybox}{关键洞察}
对解 $v=e^{-p}u$，$\partial_p v=-v$ 是恒等式，故 $H_1$ 项写成 $-H_1\partial_p v$ 不改变解。
若把 $H_2$ 项也代入（$-iH_2v=iH_2\partial_p v$），则 Fourier 后得到 $H_{\mathrm{eff}}=\xi A$，
一般\textbf{不厄米} —— 这就是"朴素提升"失败的根源。
\end{keybox}
\begin{warnbox}{正确选择}
把 $A$ 的\textbf{反厄米部分（酉生成元）保留为 Hamiltonian 的组成部分}、
只把厄米部分（耗散/增长）与 $p$ 耦合，得到厄米的 $H_{\mathrm{eff}}=\xi H_1+H_2$。
验证：$e^{-p}u(t)$ 仍满足提升方程（代入 $\partial_p v=-v$ 两边相等）。
\end{warnbox}
\begin{exbox}{三种情形的处理}
\begin{itemize}
\item 非厄米部分：吸收进 $H_2$（或厄米膨胀 $\begin{pmatrix}0&A\\A^\dagger&0\end{pmatrix}$）；
\item 非齐次项 $f(t)$：Duhamel 原理化为积分，LCU 合成；
\item 非自治 $A(t)$：$t$ 视为额外维度，化为高维自治问题（§4.6）。
\end{itemize}
\end{exbox}
\smallcite{Model 4 笔记 §4.3。}
\end{frame}

\begin{frame}{恢复定理}
\begin{Thm}[\textbf{恢复定理}]
设提升解 $w$ 满足 $\partial_t w=-H_1\partial_p w-iH_2 w$，初值 $w(0,p)=\psi(p)u_0$，
$\psi(p)=e^{-p}$ 于 $(p_*,p^*)$。令
\[
\lambda_+=\max\{\lambda_{\max}(H_1),0\},\qquad
\lambda_-=\max\{-\lambda_{\min}(H_1),0\},
\]
则对任意 $p\in I_T:=(p_*+\lambda_+T,\ p^*-\lambda_-T)$ 都有
\[
\boxed{\ u(T)=e^{p}\,w(T,p)\ }
\]
（$\lambda_\pm$ 分别是 $H_1$ 最大正特征值与最大负特征值的绝对值）。
\end{Thm}
\begin{keybox}{可行性}
$I_T$ 非空要求 $p^*-p_*>(\lambda_++\lambda_-)T$：初值"精确区间"必须足够宽。
对 $\psi=e^{-p}$（$p_*=0,p^*=\infty$）与 $H_1\preceq0$（热方程，$\lambda_+=0$）：
$I_T=(0,\infty)$，任意 $p>0$ 可恢复；$H_1$ 有正特征值时须取 $p>\lambda_+T$。
\end{keybox}
\smallcite{Model 4 笔记 §4.4。}
\end{frame}

\begin{frame}{恢复定理的证明：有限传播速度}
\begin{keybox}{第一步：参考解与误差方程}
$\tilde w(t,p):=e^{-p}u(t)$ 满足同一提升方程；误差 $e=w-\tilde w$ 满足同型方程
$\partial_t e=-H_1\partial_p e-iH_2e$，且 $e(0,p)=0$ 于 $(p_*,p^*)$。
问题归结为：\emph{初始在 $p$-区间为零的解，时刻 $T$ 在哪个区间仍为零？}
\end{keybox}
\begin{keybox}{第二步：化为纯传输方程}
$H_2$ 项 $-iH_2e$ 是 $p$-局部的（不移动 $p$-支撑）但混合 $x$-分量；
用酉变换去掉：$z=e^{iH_2t}e$ 满足
\[
\partial_t z=-e^{iH_2t}H_1e^{-iH_2t}\,\partial_p z=:-A(t)\partial_p z,
\]
其中 $A(t)$ 与 $H_1$ \textbf{酉等价}、谱相同 $\subseteq[-\lambda_-,\lambda_+]$。
传输方程的所有 $p$-信息传播速度落在 $[-\lambda_-,\lambda_+]$ 内。
\end{keybox}
\begin{keybox}{第三步：恢复区间}
初始零区域 $(p_*,p^*)$ 两侧以速度 $\lambda_+$、$\lambda_-$ 收缩，时刻 $T$ 后仍为零的区域为
$I_T=(p_*+\lambda_+T,\ p^*-\lambda_-T)$，故 $e(T,p)=0$ 于 $I_T$，即 $u(T)=e^pw(T,p)$。
\end{keybox}
\smallcite{Model 4 笔记 §4.4 证明。}
\end{frame}

\begin{frame}{恢复方式与成功概率}
\begin{keybox}{点恢复与积分恢复}
$w(T,q)=e^{-q}u(T)$ 对 $q\ge p$ 成立时，\textbf{积分恢复}更稳定：
\[
u(T)=e^p\int_p^\infty w(T,q)\,dq.
\]
\end{keybox}
\begin{keybox}{离散后的成功概率}
测量 $p$-寄存器落在恢复指标集 $I^\diamond$ 的概率近似为
\[
P_{\mathrm{suc}}\approx\frac{C_{\mathrm{rec}}^2\,\norm{u(T)}^2}{C_\psi^2\,\norm{u_0}^2},
\qquad
C_\psi^2=\sum_k|\psi(p_k)|^2,\quad
C_{\mathrm{rec}}^2=\sum_{k\in I^\diamond}e^{-2p_k}.
\]
振幅放大因子为 $O(\Gamma_\psi\,q)$，其中 $\Gamma_\psi=C_\psi/C_{\mathrm{rec}}$，
$q=\norm{u_0}/\norm{u(T)}$ 是原始非酉演化的\textbf{衰减因子}：系统强烈耗散时
（$\norm{u(T)}$ 很小），从酉嵌入中恢复目标态的概率也小 —— 这是恢复步骤的固有代价。
\end{keybox}
\smallcite{Model 4 笔记 §4.4 注。}
\end{frame}

\begin{frame}{量子模拟 PDE 的标准流程（五步）}
\begin{Def}[\textbf{五步流程}]
\begin{enumerate}
\item \textbf{薛定谔化}：warped phase transformation 把 $u_t=Lu$ 化为 $i\partial_t v=H_{\mathrm{eff}}v$（核心）；
\item \textbf{空间离散化}：对 $x$ 与辅助变量 $\xi$ 做有限差分/谱离散（$\xi$ 的截断与网格由精度决定）；
\item \textbf{Block-encoding}：把离散后的 $H_{\mathrm{eff}}$ 写成移位算子与对角算子之和，构造 $(\alpha,m)$-BE；
\item \textbf{哈密顿量模拟}：Trotter 或 qubitization/QSVT 制备 $\ket{\Psi(T)}=e^{-iH_{\mathrm{eff}}T}\ket{\Psi(0)}$；
\item \textbf{测量读出}：对 $p$（或 $\xi$）方向做逆 Fourier 变换，期望值测量提取 $u(x,T)$
（需 $O(1/\eps^2)$ 次重复）。
\end{enumerate}
\end{Def}
\begin{center}
\begin{tikzpicture}[node distance=4mm, every node/.style={font=\scriptsize,align=center}]
\node[draw,rounded corners,fill=softblue,minimum width=1.8cm,minimum height=0.7cm] (a) {薛定谔化};
\node[draw,rounded corners,fill=softgreen,minimum width=1.8cm,minimum height=0.7cm,right=of a] (b) {离散化};
\node[draw,rounded corners,fill=softyellow,minimum width=2.0cm,minimum height=0.7cm,right=of b] (c) {Block-encoding};
\node[draw,rounded corners,fill=softred,minimum width=2.2cm,minimum height=0.7cm,right=of c] (d) {Hamiltonian 模拟};
\node[draw,rounded corners,fill=white,minimum width=1.8cm,minimum height=0.7cm,right=of d] (e) {读出 $u(x,T)$};
\draw[-Latex,thick] (a)--(b); \draw[-Latex,thick] (b)--(c); \draw[-Latex,thick] (c)--(d); \draw[-Latex,thick] (d)--(e);
\end{tikzpicture}
\end{center}
\smallcite{Model 4 笔记 §4.5。}
\end{frame}

\begin{frame}{例子：热方程与对流方程}
\begin{exbox}{一维热方程（$u_t=u_{xx}$，Dirichlet 边界）}
中心差分（$h=1/(N+1)$）给出 $u_t=-L_h u$，$L_h$ 为三对角阵，特征值
$\lambda_k(L_h)=\frac{4}{h^2}\sin^2\frac{k\pi}{2(N+1)}$。
薛定谔化哈密顿量 $H_{\mathrm{eff}}(\xi)=-\xi L_h$ 对每个 $\xi$ 厄米；
$L_h$ 的三对角结构（或移位算子分解）给出 block-encoding，随后 Trotter/QSVT 演化 + 逆 Fourier 读出。
\end{exbox}
\begin{exbox}{一维对流方程（$u_t+cu_x=0$，迎风格式）}
$du_j/dt=-\frac ch(u_j-u_{j-1})$ 给出 $u_t=Au$，$A=-\frac ch(I-S_-)$（$S_-$ 向后移位）。
$A$ 非厄米；分解 $A=S+K$（$S=S^\dagger$ 厄米、$K=-K^\dagger$ 反厄米），按 $A=H_1-iH_2$ 取
$H_1=S$、$H_2=iK$：
\[
H_{\mathrm{eff}}(\xi)=\xi S+iK,
\]
对每个 $\xi$ 厄米，可直接模拟；移位算子 $S$ 的 block-encoding 由循环移位电路实现。
\end{exbox}
\begin{keybox}{要点}
两种离散矩阵（三对角/移位）都自然获得 block-encoding —— 这正是薛定谔化管线
"移位 + 对角"结构的体现。
\end{keybox}
\smallcite{Model 4 笔记 §4.6--4.7。}
\end{frame}

\begin{frame}{Smooth initialization：正则性瓶颈}
\begin{warnbox}{初值光滑性是精度的主要瓶颈}
若直接取 $\psi(p)=e^{-|p|}$，在 $p=0$ 处\textbf{不可微}。$p$ 方向用 Fourier 谱离散时，
函数越光滑谱逼近越快；对 $e^{-|p|}$ 一般需要 $\Delta p=O(\eps)$，故 $\mu_{\max}\sim1/\Delta p=O(1/\eps)$
—— 总查询复杂度对 $\eps$ 呈多项式依赖，瓶颈不在 Hamiltonian simulation，而在辅助变量初值的低光滑性。
\end{warnbox}
\begin{keybox}{Smooth initialization 的三个条件}
设 $\psi(p)=\varphi(p)e^{-p}$（$\varphi$ 为 $0\to1$ 平滑阶跃）：
\begin{enumerate}
\item \textbf{截断兼容}：$\psi$ 在计算区间外快速衰减（可周期近似）；
\item \textbf{恢复兼容}：恢复区域 $p\ge p^*$ 上 $|\psi(p)-e^{-p}|\le\eps$；
\item \textbf{导数增长控制}：对 $r\simeq\log(1/\eps)$，$\norm{\psi^{(r)}}_{L^2}^{1/r}\le Cr$
（Gevrey-1/解析型增长）。
\end{enumerate}
条件 ③ 控制谱衰减率：$r$ 阶可微 $\Rightarrow$ $\hat\psi(\xi)=O(|\xi|^{-r})$，取 $r\sim\log(1/\eps)$
得 $\mu_{\max}=O(\log(1/\eps))$ —— 光滑性以对数代价把 $\mu_{\max}$ 从 $O(1/\eps)$ 降到 $O(\log(1/\eps))$。
\end{keybox}
\smallcite{Model 4 笔记 §4.8。}
\end{frame}

\begin{frame}{Error-function 初始化与最优精度}
\begin{keybox}{Error-function 初始化}
取
\[
\varphi(p)=\frac{\mathrm{erf}(ap)+1}{2},\qquad \psi(p)=\varphi(p)e^{-p},
\]
$a\sim\sqrt{\log(1/\eps)}$、$p^*\lesssim1/2$。由 $\mathrm{erfc}(x)\lesssim e^{-x^2}$ 得恢复误差
$O(\eps)$；$\mathrm{erf}^{(k)}(ap)\sim a^kH_{k-1}(ap)e^{-a^2p^2}$ 给出解析型导数控制，从而
\[
\mu_{\max}=O\!\Big(\log\frac{1}{\eps}\Big).
\]
\end{keybox}
\begin{keybox}{矩阵查询复杂度（一次 Hamiltonian simulation）}
若 $H_1,H_2$ 的 block-encoding 归一化尺度至多为 $\alpha_A$：
\[
\widetilde O\!\Big(\Gamma_\psi\, q\,\alpha_A\,T\,\mu_{\max}\Big)
=\widetilde O\!\Big(\Gamma_\psi\,q\,\alpha_A\,T\,\log\frac{1}{\eps}\Big),
\]
在下界 $\Omega(\alpha_AT)$ 与 $\Omega(\log(1/\eps))$ 的意义下达到 matrix-query 最优
（$\Gamma_\psi=C_\psi/C_{\mathrm{rec}}$ 为恢复常数，$q=\norm{u_0}/\norm{u(T)}$ 为衰减因子）。
\end{keybox}
\smallcite{Model 4 笔记 §4.8。}
\end{frame}

\begin{frame}{带源项与时间依赖系统}
\begin{keybox}{带源项：增广齐次化}
对 $u_t=A(t)u+b(t)$，引入常向量 $r(t)$（$\dot r=0$）与 $B(t)$ 使 $B(t)r_0=b(t)$，令
$u_f=(u,r)^\top$、$A_f(t)=\begin{smallmatrix}A(t)&B(t)\\0&0\end{smallmatrix}$，则 $u_f$ 满足齐次
$\dot u_f=A_fu_f$。取 $\gamma_i=T\sup_{[0,T]}|b_i(t)|$、$B=\mathrm{diag}(b_i/\gamma_i)$ 使
$\norm{B}\le O(1/T)$。源项使扩展后的 $H_1$ 不再半负定：由 \textbf{Weyl 扰动定理}，
$\lambda_{\max}(\tilde H_1)\le\lambda_+^{\max}(H_1^A)+\delta_B$（$\delta_B=\norm{B}/2$），
恢复点取 $p^\diamond=(\lambda_+^{\max}(H_1^A)+\delta_B)T$。
\end{keybox}
\begin{keybox}{非自治系统：autonomous technique}
薛定谔化后 $\dot w=-iH(t)w$（$[H(t_1),H(t_2)]\ne0$）。引入额外变量 $s$：
$\partial_t v+\partial_s v=-iH(s)v$，$v(0,s)=G(s)w_0$；沿特征线 $s-t=\mathrm{const}$ 得
$v(t,s)=G(s-t)U_{s,s-t}w_0$，$G=\delta(s)$ 时 $w(t)=\int v(t,s)ds$。
用动量算子 $P_s=-i\partial_s$ 写成时间无关哈密顿量 $H=P_s+H(s)$；与薛定谔化结合后：
$H_h=P_s\otimes I+\sum_\ell\ket\ell\bra\ell\otimes\big(D_\mu\otimes H_1(s_\ell)-I\otimes H_2(s_\ell)\big)$。
一句话：\textbf{$p$ 变量处理非酉动力学，$s$ 变量处理时间依赖}。
\end{keybox}
\smallcite{Model 4 笔记 §4.9。}
\end{frame}

\begin{frame}{复杂度与量子优势}
\begin{Thm}[\textbf{薛定谔化的复杂度（离散化后，Hu--Jin--Liu--Zhang 2024）}]
设空间网格 $h$、维数 $d$、模拟时间 $T$、精度 $\eps$：
\begin{enumerate}
\item \textbf{热方程}（$u_t=\Delta u$）：门复杂度
$\widetilde O\big(\frac{d^2T^2\norm{u(0)}^3}{h^4\eps^3\norm{u(T)}^3}\big)$，$d>5$ 时优于经典；
\item \textbf{对流方程}（迎风）：$\widetilde O\big(\frac{dT^2\sum_\alpha c_\alpha^2\norm{u(0)}^3}{h^2\eps^3\norm{u(T)}^3}\big)$，$d>4$ 时优于经典。
\end{enumerate}
经典有限差分在 $d$ 维上代价随 $d$ \textbf{指数}增长（维度灾难），薛定谔化的门复杂度对 $d$ 是多项式 ——
\textbf{高维问题是量子优势的来源}。
\end{Thm}
\begin{warnbox}{与 HHL 路线及经典方法的对比}
隐式欧拉 + HHL 需 $M=T/\Delta t$ 次线性方程组求解（每次 QPE 与误差控制）；
薛定谔化把整个 $[0,T]$ 合并为\textbf{一次}酉演化，消除时间离散误差，且不要求离散算子厄米（迎风可直接处理）。
经典方法（有限差分/谱方法）总受 $N=O(h^{-d})$ 维度灾难限制。
前提（与 HHL 相同）：初态可高效制备、$H_{\mathrm{eff}}$ 有 block-encoding、只提取低维观测。
\end{warnbox}
\smallcite{Model 4 笔记 §4.10；Hu--Jin--Liu--Zhang, Quantum 8, 1563 (2024).}
\end{frame}

\begin{frame}{病态 PDE：薛定谔化稳定了什么}
\begin{warnbox}{病态问题为什么难：反向热方程}
$u_t=-u_{xx}$（$x\in\mathbb T$）的解满足（$\eta$ 为\textbf{空间}频率，区别于 $p$ 的频率 $\xi$）：
$\hat u(t,\eta)=e^{\eta^2(T-t)}\hat u_T(\eta)$ —— 高频模态指数增长，终值问题\textbf{不适定}。
不能要求量子算法"精确"恢复 $u(t)$，必须在频率截断意义下定义目标解并做\textbf{正则化}。
\end{warnbox}
\begin{keybox}{薛定谔化稳定的是高维系统，而非原问题}
提升后的 $w(t,x,p)$ 满足\textbf{酉}演化，高维系统本身计算稳定；但原变量 $u(t,x)$ 的恢复
仍需要正则化/截断。\textbf{薛定谔化把"数值不稳定性"与"数学不适定性"分离开来}：
前者被酉化消除，后者由用户通过频率截断显式控制。
\end{keybox}
\begin{columns}[T,totalwidth=\textwidth]
\column{0.55\textwidth}
\begin{keybox}{PDE-first vs Discretize-first}
\textbf{PDE-first}：连续层面薛定谔化 + 频率截断（$|\eta|\le\eta_{\max}$）再离散 ——
正则化误差与离散误差\textbf{分离}；
\textbf{Discretize-first}：先离散再薛定谔化 —— 恢复区被 $\lambda_{\max}(A_h)\sim h^{-2}$ 控制，
网格最高频"绑架"复杂度。\textbf{对病态 PDE 应取 PDE-first。}
\end{keybox}
\column{0.43\textwidth}
\begin{exbox}{误差分解}
\begin{center}\small
$\norm{u-u_{h,p}}\le\underbrace{\norm{u-u_{\eta_{\max}}}}_{\text{正则化}}
+\underbrace{\norm{u_{\eta_{\max}}-u_{\eta_{\max},h}}}_{\text{空间离散}}$\\[2pt]
$\qquad+\underbrace{\norm{w-w_{p,N_p}}}_{p\text{ 截断}}
+\underbrace{\text{量子误差}}_{\text{线路与采样}}$，
\end{center}
四个来源按预算统一分配。
\end{exbox}
\end{columns}
\smallcite{Model 4 笔记 §4.11；Jin--Liu--Ma, SIAM J. Sci. Comput. 2025.}
\end{frame}

% ============================================================================
\section{酉化方法族}
\begin{frame}[plain]
\centering\vspace{2.0cm}
{\Huge\bfseries\color{navy} 酉化方法族}\par\vspace{4mm}
{\large Unitarization: Schrödingerisation, LCHS, Transmutation, Dilation}
\end{frame}

\begin{frame}{统一模板与四种方法}
薛定谔化不是处理非酉动力学的唯一方法。从"把非酉传播子放进更大酉结构"的抽象视角：
\[
M(T)=e^{TA}\approx\big(\bra{\phi_{\mathrm{out}}}\otimes I\big)\,U(T)\,\big(\ket{\phi_{\mathrm{in}}}\otimes I\big).
\]
\begin{center}
\footnotesize
\begin{tabular}{@{}p{2.6cm}p{3.6cm}p{2.6cm}p{2.2cm}@{}}
\toprule
方法 & 核心变换 & 最终量子原语 & 最适合的问题\\
\midrule
薛定谔化 & 加辅助变量 $p$，$w_t=-H_1w_p+iH_2w$ & 扩大空间上的 Hamiltonian simulation & 一般线性非酉 ODE/PDE\\
LCHS & 对 $w$ 连续 Fourier 后离散 $\xi$ 求积 & 酉演化的 LCU & 积分表示清楚的系统\\
Transmutation & $A=G^\dagger G$，热半群写成波传播子 Gaussian 平均 & LCU + Hermitian dilation & 热方程、扩散型自伴问题\\
Moment-matching dilation & 构造辅助动力学 $F$ 使 $\bra\ell F^q\ket r=1$ & 大空间 HS + projection & 统一理解各酉化方法\\
\bottomrule
\end{tabular}
\end{center}
\smallcite{Model 4 笔记 §5.1。}
\end{frame}

\begin{frame}{四种方法的联系}
\begin{keybox}{LCHS $\leftarrow$ 薛定谔化}
由提升方程先连续 Fourier 再对 $\xi$ 求积：
$w(T,p)=\int\hat\psi(\xi)e^{-i\xi p}U_\xi(T)u_0\,d\xi$；截断求积后即 LCU。
取 $\psi(p)=e^{-|p|}$ 时 $\hat\psi(\xi)\propto1/(1+\xi^2)$（Cauchy 型核）。
\end{keybox}
\begin{keybox}{Transmutation}
利用 $e^{-Tx^2}=\frac{1}{2\sqrt{\pi T}}\int_\R e^{-s^2/(4T)}e^{-isx}ds$ 与
$D(G)^2=\mathrm{diag}(GG^\dagger,G^\dagger G)$，把 $e^{-TA}$ 写成酉波传播子
$e^{-isD(G)}$ 的 Gaussian 平均；适合 $A=G^\dagger G$（一阶因子 $\norm{G}=O(h^{-1})$
而非 $\norm{A}=O(h^{-2})$，相干传播代价常更优）。
\end{keybox}
\begin{keybox}{Moment-matching dilation 统一视角}
薛定谔化可看成 $F=-\partial_p$、$\ket r=e^{-p}$、$\bra\ell=\delta(p=0)$ 的
moment-matching dilation（$F^qr=r$，$\bra\ell F^q\ket r=1$）。
四种方法统一为"稳定非酉流 = 更大空间酉流的投影/平均"。
\end{keybox}
\begin{warnbox}{如何选择}
一般线性非酉系统优先薛定谔化（框架最统一、不要求特殊结构）；
已知传播子积分表示用 LCHS；正自伴扩散用 Transmutation；比较/设计新方法用 dilation 语言。
\end{warnbox}
\smallcite{Model 4 笔记 §5.2。}
\end{frame}

% ============================================================================
\section{总结与文献}
\begin{frame}[plain]
\centering\vspace{2.0cm}
{\Huge\bfseries\color{navy} 总结与文献}\par\vspace{4mm}
{\large Summary and References}
\end{frame}

\begin{frame}{本课总结：从语言到 PDE 量子求解}
\begin{enumerate}
\item \textbf{HHL}（线性代数主线）：QPE + 受控旋转 + 逆 QPE，$O(\kappa^2/\eps)$；
QSVT-HHL 以 $\log(1/\eps)$ 代替 $1/\eps$、免厄米假设；对 $N$ 指数改进但有三个"诚实前提"；
\item \textbf{哈密顿量模拟}（共同构件）：Trotter（交换子界、高阶、随机化）、Taylor/LCU、
Qubitization/QSVT（最优 $O(\alpha t+\frac{\log(1/\eps)}{\log\log(1/\eps)})$）、相互作用绘景；
无快速前推下界 $\Omega(\alpha t)$ 框定所有方法；
\item \textbf{薛定谔化}（核心）：warped phase transformation 把 $u_t=Au$（$A=H_1-iH_2$）化为
$i\partial_t\hat v=(\xi H_1+H_2)\hat v$；恢复定理（有限传播速度）；
smooth initialization 把 $\mu_{\max}$ 从 $O(1/\eps)$ 降到 $O(\log(1/\eps))$；
五步流程、源项/时间依赖、$d>4\text{--}5$ 的量子优势、病态 PDE 的 PDE-first 原则；
\item \textbf{酉化方法族}：薛定谔化 / LCHS / Transmutation / Dilation 的统一模板与选择。
\end{enumerate}
\begin{keybox}{一句话}
薛定谔化把"非酉动力学"这一 PDE 量子模拟的根本困难，
转化为"高维酉演化 + 投影恢复"，从而把 Model 1--3 的全部构件
（block-encoding、Hamiltonian simulation、振幅放大）应用到科学计算。
\end{keybox}
\end{frame}

\begin{frame}{主要文献与延伸}
\scriptsize
\begin{thebibliography}{9}
\bibitem{notes4} 牛晓铭, \emph{Model 4 薛定谔化方法}（Schrodingerization 笔记），Github 项目 Quantum-Computation.
\bibitem{day18} 湘潭大学暑期学校, \emph{偏微分方程的量子算法} day18（薛定谔化）、day19（酉化方法）课件.
\bibitem{JLY1} S. Jin, N. Liu, Y. Yu, \emph{Quantum simulation of partial differential equations via Schrödingerisation}, PRL 133, 230602 (2024); arXiv:2212.13969.
\bibitem{JLY2} S. Jin, N. Liu, Y. Yu, \emph{...via Schrödingerisation: technical details}, PRA 108, 032603 (2023); arXiv:2212.14703.
\bibitem{HJLYZ} J. Hu, S. Jin, N. Liu, L. Zhang, \emph{Quantum circuits for PDEs via Schrödingerisation}, Quantum 8, 1563 (2024); arXiv:2403.10032.
\bibitem{JLM} S. Jin, N. Liu, C. Ma, \emph{Schrödingerization based computationally stable algorithms for ill-posed problems in PDEs}, SIAM J. Sci. Comput. 47(2), B976--B1003 (2025).
\bibitem{ALL} D. An, J.-P. Liu, L. Lin, \emph{LCHS for non-unitary dynamics}, PRL 131, 150603 (2023).
\bibitem{ACL} D. An, A. Childs, L. Lin, \emph{Quantum algorithm for linear non-unitary dynamics}, Commun. Math. Phys. 407, 19 (2026).
\bibitem{Li} X. Li, \emph{From linear differential equations to unitaries: moment-matching dilation}, arXiv:2507.10285.
\bibitem{JML} S. Jin, C. Ma, E. Zuazua, \emph{Transmutation based quantum simulation}, arXiv:2601.03616.
\bibitem{HHL} A. Harrow, A. Hassidim, S. Lloyd, PRL 103, 150502 (2009); Childs--Kothari--Somma, SIAM J. Comput. 46, 1920 (2017).
\bibitem{BCCKS} Berry--Childs--Cleve--Kothari--Somma, PRL 114, 090502 (2015); Low--Chuang, Quantum 3, 163 (2019).
\end{thebibliography}
\vspace{1mm}
{\tiny\color{graytext}延伸：非线性 PDE（arXiv:2406.15821；HAM-Schrödingerisation, AAMM 2026）、Robin 边界（arXiv:2506.20478）、
Fourier 空间薛定谔化（arXiv:2505.16895）、量子线性系统的薛定谔化（arXiv:2508.13510）。
课后建议：① 手推热方程型 warped phase 与一般情形 $i\partial_t\hat v=(\xi H_1+H_2)\hat v$；
② 用特征线重证恢复定理；③ 验证 erf 初始化的 $\mu_{\max}=O(\log(1/\eps))$；④ 比较 Trotter/Taylor/QSVT 对 $H_{\mathrm{eff}}$ 的适用性。}
\end{frame}

\end{document}
